\chapter{파이썬 스크립트}

\CodeSetup{language=powershell}
\MacroIndexDisable
\TermsSetup{index=false}

\begin{flushright}
\url{https://github.com/YiHoze/texwrapper}
\end{flushright}

\begin{code}
C:\>python c:\home\bin\foo.py optional_argument mandatory_argument
C:\>foo opt_arg man_arg
\end{code}

위 줄이 아니라 아래 줄과 같은 방법으로 윈도우에서 파이썬 스크립트를 실행할 수 있게 하려면 다음과 같이 설정해야 한다.

\begin{itemize}

\item \term{PATHEXT} 환경 변수에 .py를 등록한다. .py 파일들이 실행 가능한 것들로 인식된다.

\item .py 파일들을 python.exe 또는 py.exe\annotate{python launcher}와 연결한다.

\begin{code}
C:\>cmd /c assoc .py
C:\>cmd /c ftype Python.File="C:\Users\...\python.exe" %1 %*
\end{code}

\item .py 파일들을 \term{PATH} 환경 변수에 포함된 경로에 둔다.

\end{itemize}

레이텍 사용을 비롯하여 문서 제작에 수반되는 여러 일들을 수월하게 하고자 다음과 같은 여러 래퍼\annotate{wrapper} 스크립트가 파이썬으로 작성되었다.

\begin{macros}
\item[ltx.py] lualatex 또는 xelatex을 이용하여 텍 파일을 컴파일한다.
\item[i.py] 현재 디렉토리에서 텍 파일을 찾아 ltx.py를 이용하여 컴파일한다.
\item[iu.py] 다양한 포맷의 이미지 파일들을 다른 포맷으로 변환한다.
\item[fontinfo.py] 설치된 폰트들 또는 특정 폰트의 정보를 보여준다.
\item[tlconf.py] 텍 라이브를 위한 환경을 설정한다.
\item[open.py] 텍스트 파일과 PDF 파일을 미리 지정된 프로그램으로 연다.
\item[wordig.py] 특정 문자열을 찾아 다른 문자열로 바꾼다.
\item[unit.py] 특정 단위의 값을 다른 단위로 변환한다.
\end{macros}

이들 스크립트의 대부분이 공통으로 참조하는 설정 파일인 \term{docenv.ini}는 다음과 같은 옵션들을 포함하고 있다.

\begin{codewrite}
[LaTeX]
compiler = lualatex.exe

[TeX Live]
texmflocal = C:\texlive\texmf-local\tex\latex\local

[SumatraPDF]
path = C:\Program Files\SumatraPDF\SumatraPDF.exe
\end{codewrite}
\coderead[language=ini]

\section{ltx.py: 텍 파일 컴파일하기}

\term{ltx.py}는 \term{latexmk.exe}와 비슷한 레이텍 래퍼\annotate{wrapper}이다.
텍스트 에디터에서 단축키를 이용하여 텍 파일을 컴파일할 수 있지만, \macro*{-shell-escape} 옵션을 주거나 makeindex 또는 bibtex을 돌려야 하는 것과 같은 다양한 경우들의 각각에 매우 제한적으로 서로 다른 단축키를 할당할 수 있을 뿐이다.
명령행 사용이 거북하지 않더라도, 이렇게 명령하는 것은 너무 번거롭다.

\begin{code}
C:\>latexmk -lualatex -shell-escape foo.tex
\end{code}

latexmk를 사용하지 않은 가장 큰 이유는 초보자 시절에 그것의 사용법이 거추장스러워 보였고 그래서 자연스럽게 그것에서 멀어졌다.

\begin{code}
C:\>ltx -l -s foo.tex
\end{code}

\begin{macros}
\item[-l] \term{lualatex}이 실행된다.
\item[-x] \term{xelatex}이 실행된다.
\item[-b] 이것은 \macro*{-interaction=batchmode}를 의미한다. 이 옵션이 주어지지 않으면 \macro*{-synctex=1}이 사용된다.
\item[-s] 이것은 \macro*{-shell-escape}과 같다.
\item[-w] 주어진 텍 파일이 두 번 컴파일된다.
\item[-f] 세 번 컴파일된다. .idx 파일이 있으면 texindy 또는 komkindex가 실행된다.
\item[-v] 생성된 PDF 파일이 열린다. 
\item[-n] 컴파일이 생략된다. texindy 같은 다른 프로그램만 실행하려 할 때 이 옵션이 유용하다.
\item[-i] \term{texindy}가 실행된다.
\item[-L] 색인 정렬을 위한 언어를, \macro*{german} 또는 \macro*{ger}와 같이, 지정하라. 디폴트는 \macro*{korean}이다.
\item[-k] texindy 대신 \term{komkindex}가 사용된다.
\item[-I] 색인 스타일 파일, .ist 또는 .xdy를 지정하라.
\item[-a] \macro*{-f} 옵션이 주어지면 컴파일 뒤에 .aux를 비롯한 모든 보조 파일들이 삭제된다. 그 파일들을 남겨두려면 \macro*{-a} 옵션을 사용하라.
\item[-m] 모든 색인 항목들이 PDF 북마크에 포함된다.
\item[-M] 이것은 파이썬 \term{독스트링}\annotate{docstring}을 위한 것이다. 색인 항목들 가운데 파이썬 함수들이 PDF 북마크에 포함된다.
\item[-c] 모든 보조 파일들이 삭제된다.
\item[-B] \term{bibtex}이 실행된다.
\item[-P] \term{pythontex}이 실행된다.
\end{macros}

\section{i.py: 더 게으른 사용자를 위한 레이텍 래퍼}

\term{i.py}는 ltx.py를 사용하는 레이텍 래퍼이다. 이 스크립트는 다음과 같이 작동한다.

\begin{enumerate}

\item \verb{C:\>i} \\
현재 디렉토리에서 \term{i.ini} 파일을 찾는다. 
없다면 .tex 파일들을 찾는다.
발견된 파일이 하나이면 그 파일의 이름을 아래와 같이 i.ini 파일에 기록하고 컴파일한다.

\begin{codewrite}
[tex]
target = foo.tex
\end{codewrite}
\coderead[language=ini]

\item 둘 이상의 텍 파일들이 있다면, 번호와 함께 파일 이름들이 나열된다.
번호를 입력하여 컴파일할 파일을 선택하라.
\macro*{2 5 8-10}과 같이 복수의 파일들을 선택할 수 있다.
선택된 파일들의 이름이 i.ini에 기록되고 컴파일된다.
0을 입력하면 모든 파일들이 i.ini에 기록되고 컴파일된다.

\item i.ini 파일에 기록된 텍 파일이 하나이면 그것이 컴파일된다.
둘 이상이면, 사용자가 선택할 수 있도록 번호와 함께 파일 이름들이 나열된다. 
\macro*{2 5 8-10}과 같이 복수의 파일들을 선택할 수 있다.
0을 입력하면 모든 파일들이 컴파일된다.

\item \verb{C:\>i foo} \\
현재 디렉토리에 i.ini 파일이 없다면, 파일 이름에 \macro*{foo}를 포함하는 텍 파일들을 찾는다.
나머지 과정은 첫째 및 둘째 단계와 동일한다.

\item i.ini 파일이 있고, 기록된 파일 이름들 가운데 \macro*{foo}를 포함하는 것이 하나이면 그것이 컴파일된다.
나머지 과정은 셋째 단계와 동일하다.

\end{enumerate}

i.ini는, 다양한 방법으로 처리할 수 있도록, 여러 옵션들을 갖고 있다.

\begin{codewrite}
[tex]
target = foo.tex
compiler = 
draft = wordig.py -a "..." -s "..." %(target)s
final = wordig.py -a "..." -s "..." %(target)s
final_compiler = 
after = 
main = \input{preamble}
    \begin{document}
    \maketitle
    \input{\1}
    \end{document}
\end{codewrite}
\coderead[language=ini]

i.py를 다음과 같은 선택 옵션들과 함께 사용할 수 있다. i.py가 인식하지 못하는 옵션들은 ltx.py에 전달된다.
i.ini의 \macro{compiler}에 설정된 것들도 ltx.py에 전달된다.

\begin{macros}
\item[-U] i.ini를 업데이트할 수 있도록 텍 파일들을 나열한다.

\item[-D] 초본을 만들기 위한 전처리로서 i.ini의 \macro{draft}에 지정된 명령이 수행된다.

\item[-F] 최종본을 만들기 위한 전처리로서 i.ini의 \macro{final}에 지정된 명령이 수행된다.
그리고 \macro{final_compiler}에 설정된 옵션들이 ltx.py에 전달된다.

\item[-W] 이것은 여러 하위 파일들을 따로 컴파일하기 위한 목적으로 고안되었다.
i.ini의 \macro{main} 옵션이 위 예와 같이 설정되어야 한다. 
\macro*{\1}이 선택된 파일로 바뀐다.

\item[-N] 컴파일이 생략된다.

\item[-C] i.ini 파일이 생성된다.
\end{macros}

\section{iu.py: 이미지를 다른 포맷으로 변환하기}

\term{iu.py}에 정의된 클래스 이름이 ImageUtility인데, 그것은 억지로 지어 붙인 것이고, 이 파일 이름은 그림처럼 이쁜 아이유를 기리기 위한 것이다.
이것 역시 래퍼이기 때문에 다음과 같은 프로그램들이 필요하다.

\begin{itemize}*
\item \term{이미지매직}\annotate{ImageMagick}
\item \term{잉크스케이프}\annotate{Inkscape}
\item \term{텍 라이브}\annotate*{epstopdf.exe, pdfcrop.exe,pdftops.exe}
\end{itemize}

iu.py는 다음과 같은 것들을 할 수 있다.

\begin{itemize}

\item 비트맵 이미지들을 다른 포맷의 이미지로 변환할 수 있다.

\begin{code}
C:\>iu -t png *.jpg
\end{code}

\item 비트맵 이미지들을 PDF 또는 EPS로 변환할 수 있다.

\begin{code}
C:\>iu -t pdf *.png
\end{code}

\item 벡터 이미지들을 비트맵 또는 다른 포맷의 벡터 이미지로 변환할 수 있다.
단 SVG를 PDF나 EPS로 변환할 수 있지만 그 반대로는 가능하지 않다.

\begin{code}
C:\>iu -t png *.eps
C:\>iu -t eps *.svg
\end{code}

\item 비트맵 이미지들의 픽셀 크기나 해상도를 변경할 수 있다.

\begin{code}
C:\>iu -r -s 75 *.jpg
C:\>iu -r -d 150 *.jpg
\end{code}

\item 비트맵 이미지들의, 픽셀 크기를 포함하는, 이미지 정보를 볼 수 있다.

\begin{code}
C:\>iu -i foo.jpg
\end{code}
\end{itemize}

\begin{macros}
\item[-t] 이 옵션과 함께 지정된 포맷\annotate*{eps, pdf, svg, bmp, cr2, gif, jpg, jpeg, pbm, png, ppm, tga, tiff, webp}으로 변환된다. 
디폴트는 \macro*{pdf}이다.
GIF 이미지와 여러 페이지로 이루어진 PDF 파일은 낱개 이미지로 쪼개진다. 

\item[-r] 비트맵 이미지의 크기를 변경하려면 이 옵션을 지정하라.

\item[-d] 해상도를 지정하라. 단위는 \term{ppc}\annotate{pixels per centimeter}이다.

\item[-m] 아래 예의 경우에 가로 800 픽셀보다 큰 이미지들이 800 픽셀로 축소된다.
\begin{code}
C:\>iu -r -m 800 *.png
\end{code}

\item[-s] \macro*{75}를 지정하면 원래 크기의 75 퍼센트로 바뀐다.

\item[-R] 모든 하위 디렉토리에 있는 이미지 파일들이 처리된다.
\end{macros}

\section{fontinfo.py: 폰트 정보 보기}

\term{fontinfo.py}는 텍 라이브에 포함된 \term{fc-list.exe}와 \term{otfinfo.exe}를 이용하여 설치된 폰트들에 대한 정보를 보여준다.
사용 방법은 다음과 같다.

\begin{enumerate}

\item 텍 라이브를 포함하여 설치된 모든 폰트들의 목록이 \macro{fonts_list.txt}에 만들어진다.
출력 파일의 이름을 바꾸려면 \macro*{-o} 옵션을 사용하라.

\begin{code}
C:\>fontinfo 
C:\>fontinfo -o myfonts.txt
\end{code}

\item 폰트 이름 또는 파일 이름을 지정하면, 언어별로 짧은 문장을 보여주는, 같은 이름의 PDF가 만들어진다. 
이 기능은 mytex.py와 latex.db를 요구한다.
\ptref{sec:mytex}\을 보라.

\begin{code}
C:\>fontinfo "Noto Serif"
C:\>fontinfo NotoSerif-Regular.ttf
\end{code}

\item 폰트 정보를 보려면 \macro*{-i} 옵션과 함께 폰트 이름 또는 파일 이름을 지정하라.

\begin{code}
C:\>fontinfo -i "Noto Serif"
C:\>fontinfo -i NotoSerif-Regular.ttf
\end{code}

\end{enumerate}

\section{tlconf.py: 텍 라이브와 관련된 설정}

텍 라이브를 처음 설치한 다음에 설정해야 할 것들이 몇 가지 있다.
\term{tlconf.py}는 텍 라이브를 위한 설정들을 수월하게 처리하고자 고안되었다.
\term{docenv.ini} 설정 파일에서 다음 옵션들이 적절하게 설정되어 있어야 한다.

\begin{codewrite}
[TeX Live]
texmflocal = C:\texlive\texmf-local\tex\latex\local
repository_main = https://cran.asia/tex/systems/texlive/tlnet/
repository_private = http://ftp.ktug.org/KTUG/texlive/tlnet

TEXEDIT = \"C:\...\Microsoft VS Code\code.exe\" -r -g \"%%s\":%%d
TEXMFHOME = C:\home\texmf

[LOCAL.CONF]
path = C:\texlive\202x\texmf-var\fonts\conf\local.conf
content = <dir>C:/Users/.../AppData/Local/Microsoft/Windows/Fonts</dir>

[SumatraPDF]
path = C:\Program Files\SumatraPDF\SumatraPDF.exe
inverse-search = "C:\...\Microsoft VS Code\code.exe" -r -g "%%f":%%l
\end{codewrite}
\coderead[language=ini] 

텍 라이브를 업그레이드한 경우에도 \term{local.cnf}는 손봐야 한다.

\begin{code}
C:\>tlconf -C 
\end{code}

\begin{macros}
\item[-H] \term{TEXMFHOME} 환경 변수가 만들어지고 설정 파일에 지정되어 있는 경로로 설정된다.
% \item[-C] \term{texmf.cnf} 파일에 사용자의 폰트 디렉토리 경로가 추가된다.\footnote{TeX Live 2021부터 이 조작은 더 이상 필요하지 않다.}
\item[-L] \term{local.cnf} 파일이 만들어지고 사용자의 폰트 디렉토리 경로가 기재된다.
\item[-e] \term{TEXTEDIT} 환경 변수가 만들어진다.
\item[-p] \term{수마트라 PDF}\annotate{Sumatra PDF}의 경로와 \term{인버스 서치}\annotate{inverse search} 옵션이 설정된다. 다른 PDF 뷰어를 지정할 수 있지만, 인버스 서치 기능을 제공하는 PDF 뷰어가 더 있는지 모르겠다.
\item[-r] \term{tlmgr.exe}가 사용할 텍 라이브 저장소가 지정된다.
\item[-u] 텍 라이브가 업데이트된다.
\item[-c] 지텍을 위한 폰트 캐시가 업데이트된다.
\item[-l] 루아텍을 위한 폰트 데이터베이스가 업데이트된다.
\item[-b] 이것은 위의 모든 옵션을 선택한 것과 동일하다.
\item[-q] 사용자 확인을 묻지 않고 진행된다.
\end{macros}

HzGuide 클래스와 함께 제공되는 \term{tlconf.cmd}는 tlconf.py의 간단 버전이다.

\begin{codewrite}
C:\>tlconf.cmd cnf
C:\>tlconf.cmd texedit
C:\>tlconf.cmd texmfhome
C:\>tlconf.cmd sumatrapdf
C:\>tlconf.cmd batch
\end{codewrite}
\powershellread

아무 인자도 주지 않으면 각 옵션에 대한 도움말이 나온다.

\begin{macros}<tlconf.cmd>
\item[cnf] 사용자 폰트 폴더의 경로를 포함하는 \macro{local.cnf}를 생성한다.
\item[texedit] \macro{TEXEDIT} 환경 변수를 설정한다.
\item[texmfhome] \macro{TEXMFHOME} 환경 변수를 설정한다.
\item[sumatrapdf] \term{Sumatra PDF}에 \term{인버스 서치}\annotate{inverse search} 옵션을 설정한다.
\item[batch] 이것은 위의 모든 옵션들을 지시한 것과 동일하다.
\end{macros}

\coderead[file=tlconf.cmd, language=bat, fontsize=\footnotesize]

\section{open.py: 파일 열기}

탐색기에서 어떤 파일을 열기 위해 특정 프로그램과 연결하는 일은 다소 성가시다.
\term{open.py}는 docenv.ini에 지정된 프로그램들을 이용하여 파일들을 연다.

\begin{codewrite}
[SumatraPDF]
path = C:\Program Files\SumatraPDF\SumatraPDF.exe

[Text Editor]
path = C:\...\Microsoft VS Code\code.exe
associations = .aux, .bib, .bst, .cls, .cnf, .conf, .cmd, .css, .csv, .gv,
 .ini, .ind, .idx, .ist, .ipynb, .list, .lof, .log, .lot, .md, .ps1, .py, 
 .rst, .sty, .tex, .tmp, .toc, .tsv, .txt

[Adobe Reader]
path = C:\Program Files (x86)\Adobe\Acrobat Reader DC\Reader\AcroRd32.exe

[Web Browser]
path = C:\Program Files (x86)\Google\Chrome\Application\chrome.exe
\end{codewrite}
\coderead[language=ini] 

아래 예처럼 대상을 다섯 가지로 나눌 수 있다.

\begin{code}
C:\>open foo.tex
C:\>open foo.pdf
C:\>open -s memoir.cls 
C:\>open -w ktug.org
C:\>open foo.docx
\end{code}

명시되지 않은 확장자의 파일을 여는 것은 윈도우가 결정한다.

\begin{macros}
\item[-a] 기본 에디터 대신 다른 프로그램을 사용하려면 이 옵션과 함께 사용할 프로그램을 지정하라.
\item[-o] 이 옵션과 함께 주어진 인자가 에디터에 전달된다.
\item[-A] PDF 파일을 여는 데에 아크로뱃 리더가 사용된다.
\item[-s] 텍 라이브에 포함된 파일을 열려면 이 옵션을 사용하라.
\item[-f] 이미지 파일 같은 바이너리 파일을 포함하여, 지정되지 않은 형식의 파일을 텍스트 에디터로 열려면 이 옵션을 사용하라.
\item[-w] 웹 주소를 열려면 이 옵션을 사용하라.
\end{macros}

\section{wordig.py: 문자열을 찾아 바꾸기}

\term{wordig.py}는, 단어 수 세기를 비롯하여, 레이텍 사용자가 텍스트 파일을 갖고 할 법한 거의 모든 기능을 제공한다.

\begin{macros}
\item[-a] 이 옵션과 함께 찾을 문자열 또는 정규 표현식을 지정하라.
\item[-s] 바꿔치기할 문자열을 지정하라.
\item[-b] 백업 파일이 만들어진다.
\item[-P] 지정된 \term{TSV} 또는 \term{CSV} 파일에 포함된 패턴들에 따라 문자열 바꾸기가 이루어진다.
\item[-J] 조사들이 \term{자동 조사} 매크로로 바뀐다. 예를 들어 \macro*{을}이 \macro*{\을}로 바뀐다.
\item[-j] 자동 조사 매크로들이 조사로 바뀐다.
\item[-e] 지정된 파일에 포함된 패턴들에 따라 문자열들이 추출되어 새 파일에 저장된다.
\item[-w] 모든 단어들이 추출된다.
\item[-t] 텍 매크로들이 추출된다.
\item[-d] 주어진 파일에서 텍 매크로들이 제거된다.
\item[-o] 출력 파일의 이름을 지정하라. 디폴트는 옵션에 따라 다르다.
\item[-r] 모든 하위 디렉토리에 있는 파일들이 처리된다.
\item[-P] 각 PDF 파일들의 페이지 수와 합계를 구하려면 이 옵션을 사용하라.
\item[-U] 주어진 문자들의 유니코드 정보가 표시된다.
\item[-c] \term{EUC-KR}로 인코딩된 파일을 \term{UTF-8}로 변환한다.
\end{macros}

아무 옵션도 지시하지 않으면 주어진 파일들의 글자 수와 단어 수가 표시된다.
PDF 파일은 .txt 파일로 변환된 다음에 처리된다.

\begin{code}
C:\>wordig *.tex
C:\>wordig *.pdf
\end{code}

문자열을 찾거나 바꾸는 방법은 다음과 같다.

\begin{code}
C:\>wordig -a itemize *.tex
C:\>wordig -a itemize -s enumerate *.tex
C:\>wordig -p foo.tsv *.tex
\end{code}

텍 매크로를 제거하는 데 사용되는 detex.tsv 파일의 일부가 다음과 같다.

\begin{codewrite}
(?<!\\)%.+?\n	\n
\\title\{(.+?)\}	\1
\\chapter\{(.+?)\}	\1
\\section\{(.+?)\}	\1
\\caption\{(.+?)\}	\1
\\textbf\{(.+?)\}	\1
\end{codewrite}
\coderead[language=text]

\section{unit.py: 단위 변환하기}

\term{unit.py}는 다음과 같은 여러 비표준 단위의 값들을 미터법 단위로 환산한다.\annotate{미국이 미터법을 채용하지 않은 거의 유일한 나라이다.}

\begin{codewrite}
acre, degree, fahrenheit, foot, inch, knot, mile, pascal, point, pound, pyeong, yard
\end{codewrite}
\coderead[language=text]

다른 단위와 구분될 수 있을 만큼 사용할 단위의 처음 글자들을 숫자와 함께 지정하라.
아래 예에서 첫 줄의 \macro*{p}는 \macro*{pascal}로, 둘째 줄의 \macro*{po}는 \macro*{point}로 간주된다.

\begin{code}
C:\>unit 10 p 
C:\>unit 20 po
C:\>unit 30 pou
C:\>unit 30 py
\end{code}

천 자리 구분자로 쉼표를 사용할 수 있는데, 파워셸에서는 따옴표로 값을 감싸거나 이스케이프 문자인 억음 부호와 함께 쉼표를 써야 한다.

\begin{code}
C:\>unit 2`,000`,000 mi
C:\>unit "2,000,000" mi
\end{code}

unit.py는 또한 RGB 색상을 CMYK로 또는 그 반대로 바꿀 수 있다.
RGB는 HTML 색상 값이나 16 비트 10진수로, CMYK는 백분율 값으로 지정해야 한다.
띄어쓰기를 하려면 값을 따옴표로 감싸야 한다.

\begin{code}
C:\>unit -c 0000FF
C:\>unit -c 240,120,99
C:\>unit -c '240, 120, 99'
C:\>unit -c "0, 0.1, 0.33, 0.01"
\end{code}

\section{레이텍 템플릿 모음}
\label{sec:mytex}

여러 레이텍 예들이 \term{latex.db}에 저장되어 있다.
\term{my.tex}이 latex.db에서 지정된 템플릿을 꺼내어 저장하고 컴파일한다.
다음은 latex.db의 일부이다.

\begin{codewrite}
[multilingual]
description = Use this template to see how many languaegs are supported by a font.
compiler = -c
tex_output = multilingual.tex
placeholders = 1
defaults = Noto Serif
tex = 
	\documentclass{minimal}
	\usepackage{fontspec}
	\setlength\parskip{1.25\baselineskip}
	\setlength\parindent{0pt}
	\setmainfont{\1}
	\begin{document}
	0 1 2 3 4 5 6 7 8 9
	
	english:
	© Keep this manual® for future use. These---lines--are-dashes.
	
	chinese 中文:
	请妥善保存本手册，以备将来使用。
\end{codewrite}
\coderead[language=ini] 

대괄호로 감싼 말이, ini 형식에서 섹션을 가리키는데, 템플릿을 나타낸다.

\begin{code}
C:\>mytex multilingual -s "Noto Serif CJK KR"
\end{code}

\term{multilingual} 템플릿은 지정된 폰트가 얼마나 많은 언어들을 지원할 수 있는지 보기 위해 고안된 것이다.

\image[float, frame, scale=0.75]{multilingual_example}(multilingual.pdf)

\begin{macros}
\item[-o] 출력 파일의 이름을 지정하라.
\item[-s] \macro*{placeholders} 옵션이 지정된 템플릿에서 \macro*{\1}처럼 백슬래시로 시작하는 숫자가 이 옵션과 함께 주어진 문자열로 대체된다.
\item[-n] \macro*{compiler}에 설정된 옵션들이 ltx.py에 전달되는데, \macro*{-n}가 주어지면 ltx.py가 호출되지 않는다.
\item[-f] \macro*{compiler} 옵션에 아무 것도 설정되어 있지 않아도 템플릿 파일이 ltx.py에 의해 컴파일된다.
\item[-r] PDF 파일을 제외하고 템플릿 파일을 비롯하여 부수적으로 생성된 모든 파일들이 컴파일 뒤에 삭제된다.
\item[-l] 템플릿들이 나열된다.
\item[-L] 각 템플릿의 설명과 함께 템플릿들이 나열된다.
\item[-d] 지정된 템플릿에 대한 상세한 설명이 표시된다.
\item[-i] 새로 작성한 텍 파일을 latex.db에 삽입한다. 텍 파일의 이름이 템플릿 이름이 된다.
\begin{code}
C:\>mytex -i foo.tex
\end{code}

\item[-u] 추출된 텍 파일을 수정하고 그것으로 latex.db를 업데이하려면 이 옵션을 사용하라.
\begin{code}
C:\>mytex -u multilingual
\end{code}

\item[-b] 모든 템플릿들이 추출된다.
\end{macros}

mytex.py가 인식하지 못하는 옵션들은 ltx.py에 전달된다

\section{이따금 요긴한 스크립트들}
\label{sec:myscript}

\term{scripts.db}에 가끔 유용할 수 있는 스크립트들이 들어있다. 
대부분이 파이썬으로 작성되었고, 일부는 파워셸로 작성되었다.
\term{myscript.py}를 이용하여 특정 스크립트를 꺼내어 실행할 수 있다.
다음 예는 scripts.db에서 gencode.py를 꺼내어 저장하면서 KTUG을 가리키는 QR 코드를 만드는 방법을 보여준다.

\begin{code}
C:\>myscript gencode -q www.ktug.org
\end{code}

\begin{macros}
\item[-I] 새로 작성한 스크립트 파일을 scripts.db에 삽입한다.
\begin{code}
C:\>myscript -I foo.py
C:\>myscript -I foo.ps1
\end{code}

\item[-U] 추출된 스크립트 파일을 수정하고 그것으로 scripts.db를 업데이트하려면 이 옵션을 사용하라.
\begin{code}
C:\>myscript -U foo
\end{code}

\item[-B] 모든 스크립트들이 추출된다.
\item[-C] 실행 뒤에 추출된 스크립트가 삭제된다.
\item[-R] 선택된 스크립트가 데이터베이스에서 제거된다.
\end{macros}
